\documentclass[UTF8,12pt]{ctexart}
\usepackage[a4paper,margin=24mm]{geometry}
\usepackage{amsmath,amssymb,bm,graphicx,adjustbox}
\newcommand{\fitmath}[1]{\begin{adjustbox}{max width=\linewidth}$\displaystyle #1$\end{adjustbox}}
\setlength{\parindent}{0pt}
\setlength{\parskip}{.7em}
\sloppy
\allowdisplaybreaks
\begin{document}
\section*{2012 年考研数学三 · 真题}
\subsection*{原 PDF 第 36 页}

\subsection*{2012 年全国硕士研究生招生考试试题}

\subsection*{一、选择题（本题共 8 小题，每小题 4 分，共 32 分。在每小题给出的四个选项中，只有一项符合题目要求，把所选项前的字母填在题后的括号内。）}

(1) 曲线 \(\displaystyle y=\dfrac{\displaystyle x^2+x}{\displaystyle x^2-1}\) 的渐近线的条数为（　）。


(A) 0。 (B) 1。 (C) 2。 (D) 3。


(2) 设函数 \(\displaystyle f(x)=(e^x-1)(e^{2x}-2)\cdots(e^{nx}-n)\)，其中 \(\displaystyle n\) 为正整数，则 \(\displaystyle f^{\prime}(0)=\)（　）。


(A) \(\displaystyle (-1)^{n-1}(n-1)!\)。 (B) \(\displaystyle (-1)^n(n-1)!\)。 (C) \(\displaystyle (-1)^{n-1}n!\)。 (D) \(\displaystyle (-1)^nn!\)。


(3) 设函数 \(\displaystyle f(t)\) 连续，则二次积分 \(\displaystyle \int_0^{\pi/2}d\theta\displaystyle \int_{2\cos\theta}^2 f(r^2)r\,dr=\)（　）。


(A) \(\displaystyle \int_0^2dx\displaystyle \int_{\sqrt{2x-x^2}}^{\sqrt{4-x^2}}\sqrt{x^2+y^2}f(x^2+y^2)\,dy\)。 (B) \(\displaystyle \int_0^2dx\displaystyle \int_{\sqrt{2x-x^2}}^{\sqrt{4-x^2}}f(x^2+y^2)\,dy\)。


(C) \(\displaystyle \int_0^2dy\displaystyle \int_{1+\sqrt{1-y^2}}^{\sqrt{4-y^2}}\sqrt{x^2+y^2}f(x^2+y^2)\,dx\)。 (D) \(\displaystyle \int_0^2dy\displaystyle \int_{1+\sqrt{1-y^2}}^{\sqrt{4-y^2}}f(x^2+y^2)\,dx\)。


(4) 已知级数 \(\displaystyle \sum_{n=1}^{\infty}(-1)^n\sqrt n\sin\dfrac{\displaystyle 1}{\displaystyle n^\alpha}\) 绝对收敛，级数 \(\displaystyle \sum_{n=1}^{\infty}\dfrac{\displaystyle (-1)^n}{\displaystyle n^{2-\alpha}}\) 条件收敛，则（　）。


(A) \(\displaystyle 0<\alpha\leq\dfrac{\displaystyle 1}{\displaystyle 2}\)。 (B) \(\displaystyle \dfrac{\displaystyle 1}{\displaystyle 2}<\alpha\leq1\)。 (C) \(\displaystyle 1<\alpha\leq\dfrac{\displaystyle 3}{\displaystyle 2}\)。 (D) \(\displaystyle \dfrac{\displaystyle 3}{\displaystyle 2}<\alpha<2\)。


(5) 设 \(\displaystyle \alpha_1=\begin{pmatrix}0\\0\\c_1\end{pmatrix},\allowbreak \ \alpha_2=\begin{pmatrix}0\\1\\c_2\end{pmatrix},\allowbreak \ \alpha_3=\begin{pmatrix}1\\-1\\c_3\end{pmatrix},\allowbreak \ \alpha_4=\begin{pmatrix}-1\\1\\c_4\end{pmatrix}\)，其中 \(\displaystyle c_1,\allowbreak c_2,\allowbreak c_3,\allowbreak c_4\) 为任意常数，则下列向量组线性相关的为（　）。


(A) \(\displaystyle \alpha_1,\allowbreak \alpha_2,\allowbreak \alpha_3\)。 (B) \(\displaystyle \alpha_1,\allowbreak \alpha_2,\allowbreak \alpha_4\)。 (C) \(\displaystyle \alpha_1,\allowbreak \alpha_3,\allowbreak \alpha_4\)。 (D) \(\displaystyle \alpha_2,\allowbreak \alpha_3,\allowbreak \alpha_4\)。


(6) 设 \(\displaystyle A\) 为 3 阶矩阵，\(\displaystyle P\) 为 3 阶可逆矩阵，且 \(\displaystyle P^{-1}AP=\begin{pmatrix}1&0&0\\0&1&0\\0&0&2\end{pmatrix}\)。若 \(\displaystyle P=(\alpha_1,\allowbreak \alpha_2,\allowbreak \alpha_3),\allowbreak \ Q=(\alpha_1+\alpha_2,\allowbreak \alpha_2,\allowbreak \alpha_3)\)，则 \(\displaystyle Q^{-1}AQ=\)（　）。


(A) \(\displaystyle \begin{pmatrix}1&0&0\\0&2&0\\0&0&1\end{pmatrix}\)。 (B) \(\displaystyle \begin{pmatrix}1&0&0\\0&1&0\\0&0&2\end{pmatrix}\)。 (C) \(\displaystyle \begin{pmatrix}2&0&0\\0&1&0\\0&0&2\end{pmatrix}\)。 (D) \(\displaystyle \begin{pmatrix}2&0&0\\0&2&0\\0&0&1\end{pmatrix}\)。


(7) 设随机变量 \(\displaystyle X\) 与 \(\displaystyle Y\) 相互独立，且都服从区间 \(\displaystyle (0,\allowbreak 1)\) 上的均匀分布，则 \(\displaystyle P\{X^2+Y^2\leq1\}=\)（　）。


(A) \(\displaystyle \dfrac{\displaystyle 1}{\displaystyle 4}\)。 (B) \(\displaystyle \dfrac{\displaystyle 1}{\displaystyle 2}\)。 (C) \(\displaystyle \dfrac{\displaystyle \pi}{\displaystyle 8}\)。 (D) \(\displaystyle \dfrac{\displaystyle \pi}{\displaystyle 4}\)。


\clearpage

\subsection*{原 PDF 第 37 页}

(8) 设 \(\displaystyle X_1,\allowbreak X_2,\allowbreak X_3,\allowbreak X_4\) 为来自总体 \(\displaystyle N(1,\allowbreak \sigma^2)\)（\(\displaystyle \sigma>0\)）的简单随机样本，则统计量 \(\displaystyle \dfrac{\displaystyle X_1-X_2}{\displaystyle |X_3+X_4-2|}\) 的分布为（　）。


(A) \(\displaystyle N(0,\allowbreak 1)\)。 (B) \(\displaystyle t(1)\)。 (C) \(\displaystyle \chi^2(1)\)。 (D) \(\displaystyle F(1,\allowbreak 1)\)。


\subsection*{二、填空题（本题共 6 小题，每小题 4 分，共 24 分，把答案填在题中横线上。）}

(9) \(\displaystyle \lim_{x\to\pi/4}(\tan x)^{1/(\cos x-\sin x)}=\)\_\_\_\_\_\_。


(10) 设函数 \(\displaystyle f(x)=\begin{cases}\ln\sqrt x,\allowbreak &x\geq1,\allowbreak \\2x-1,\allowbreak &x<1,\allowbreak \end{cases}\ y=f(f(x))\)，则 \(\displaystyle \left.\dfrac{\displaystyle dy}{\displaystyle dx}\right|_{x=e}=\)\_\_\_\_\_\_。


(11) 设连续函数 \(\displaystyle z=f(x,\allowbreak y)\) 满足 \(\displaystyle \lim_{\substack{x\to0\\y\to1}}\dfrac{\displaystyle f(x,\allowbreak y)-2x+y-2}{\displaystyle \sqrt{x^2+(y-1)^2}}=0\)，则 \(\displaystyle dz|_{(0,\allowbreak 1)}=\)\_\_\_\_\_\_。


(12) 由曲线 \(\displaystyle y=\dfrac{\displaystyle 4}{\displaystyle x}\) 和直线 \(\displaystyle y=x\) 及 \(\displaystyle y=4x\) 在第一象限中围成的平面图形的面积为\_\_\_\_\_\_。


(13) 设 \(\displaystyle A\) 为 3 阶矩阵，\(\displaystyle |A|=3\)，\(\displaystyle A^*\) 为 \(\displaystyle A\) 的伴随矩阵，若交换 \(\displaystyle A\) 的第 1 行与第 2 行得矩阵 \(\displaystyle B\)，则 \(\displaystyle |BA^*|=\)\_\_\_\_\_\_。


(14) 设 \(\displaystyle A,\allowbreak B,\allowbreak C\) 是随机事件，\(\displaystyle A\) 与 \(\displaystyle C\) 互不相容，\(\displaystyle P(AB)=\dfrac{\displaystyle 1}{\displaystyle 2},\allowbreak \ P(C)=\dfrac{\displaystyle 1}{\displaystyle 3}\)，则 \(\displaystyle P(AB\mid\overline C)=\)\_\_\_\_\_\_。


\subsection*{三、解答题（本题共 9 小题，共 94 分，解答应写出文字说明、证明过程或演算步骤。）}

(15)（本题满分 10 分）求极限 \(\displaystyle \lim_{x\to0}\dfrac{\displaystyle e^{x^2}-e^{2-2\cos x}}{\displaystyle x^4}\)。


(16)（本题满分 10 分）计算二重积分 \(\displaystyle \iint_D e^xxy\,dx\,dy\)，其中 \(\displaystyle D\) 是以曲线 \(\displaystyle y=\sqrt x,\allowbreak \ y=\dfrac{\displaystyle 1}{\displaystyle \sqrt x}\) 及 \(\displaystyle y\) 轴为边界的无界区域。


(17)（本题满分 10 分）某企业为生产甲、乙两种型号的产品投入的固定成本为 10 000（万元）。设该企业生产甲、乙两种产品的产量分别为 \(\displaystyle x\)（件）和 \(\displaystyle y\)（件），且这两种产品的边际成本分别为 \(\displaystyle 20+\dfrac{\displaystyle x}{\displaystyle 2}\)（万元／件）与 


\clearpage

\subsection*{原 PDF 第 38 页}

\(\displaystyle 6+y\)（万元／件）。（I）求生产甲、乙两种产品的总成本函数 \(\displaystyle C(x,\allowbreak y)\)（万元）；（II）当总产量为 50 件时，甲、乙两种产品的产量各为多少时可使总成本最小？求最小总成本；（III）求总产量为 50 件且总成本最小时甲产品的边际成本，并解释其经济意义。


(18)（本题满分 10 分）证明 \(\displaystyle x\ln\dfrac{\displaystyle 1+x}{\displaystyle 1-x}+\cos x\geq1+\dfrac{\displaystyle x^2}{\displaystyle 2}\)（\(\displaystyle -1<x<1\)）。


(19)（本题满分 10 分）已知函数 \(\displaystyle f(x)\) 满足方程 \(\displaystyle f^{\prime\prime}(x)+f^{\prime}(x)-2f(x)=0\) 及 \(\displaystyle f^{\prime\prime}(x)+f(x)=2e^x\)。


（I）求 \(\displaystyle f(x)\) 的表达式；（II）求曲线 \(\displaystyle y=f(x^2)\displaystyle \int_0^x f(-t^2)\,dt\) 的拐点。


(20)（本题满分 11 分）设 \(\displaystyle A=\begin{pmatrix}1&a&0&0\\0&1&a&0\\0&0&1&a\\a&0&0&1\end{pmatrix},\allowbreak \ \beta=\begin{pmatrix}1\\-1\\0\\0\end{pmatrix}\)。


（I）计算行列式 \(\displaystyle |A|\)；（II）当实数 \(\displaystyle a\) 为何值时，方程组 \(\displaystyle Ax=\beta\) 有无穷多解，并求其通解。


\clearpage

\subsection*{原 PDF 第 39 页}

(21)（本题满分 11 分）已知 \(\displaystyle A=\begin{pmatrix}1&0&1\\0&1&1\\-1&0&a\\0&a&-1\end{pmatrix}\)，二次型 \(\displaystyle f(x_1,\allowbreak x_2,\allowbreak x_3)=x^{\mathrm T}(A^{\mathrm T}A)x\) 的秩为 2。


（I）求实数 \(\displaystyle a\) 的值；（II）求正交变换 \(\displaystyle x=Qy\) 将 \(\displaystyle f\) 化为标准形。


(22)（本题满分 11 分）设二维离散型随机变量 \(\displaystyle (X,\allowbreak Y)\) 的概率分布为：\(\displaystyle \begin{array}{c|ccc}X\backslash Y&0&1&2\\\hline 0&\dfrac{\displaystyle 1}{\displaystyle 4}&0&\dfrac{\displaystyle 1}{\displaystyle 4}\\1&0&\dfrac{\displaystyle 1}{\displaystyle 3}&0\\2&\dfrac{\displaystyle 1}{\displaystyle 12}&0&\dfrac{\displaystyle 1}{\displaystyle 12}\end{array}\)。


（I）求 \(\displaystyle P\{X=2Y\}\)；（II）求 \(\displaystyle \operatorname{Cov}(X-Y,\allowbreak Y)\)。


(23)（本题满分 11 分）设随机变量 \(\displaystyle X\) 与 \(\displaystyle Y\) 相互独立，且都服从参数为 1 的指数分布。记 \(\displaystyle U=\displaystyle \max\{X,\allowbreak Y\},\allowbreak \ V=\displaystyle \min\{X,\allowbreak Y\}\)。


（I）求 \(\displaystyle V\) 的概率密度 \(\displaystyle f_V(v)\)；（II）求 \(\displaystyle E(U+V)\)。


\clearpage
\end{document}
